%Môn: Toán
\begin{ex}
    Cho cấp số cộng $(u_n)$ có $u_1 = 3$ và công sai $d = 4$. Tính số hạng $u_2$.
    \shortans{7}
    \loigiai{$u_2 = u_1 + d = 3 + 4 = 7$.}
\end{ex}
%Môn: Vật Lí
\begin{ex}
    Một vật dao động điều hòa với tần số $f = 50$ Hz. Chu kì dao động $T$ của vật là bao nhiêu giây?
    \shortans{0,02}
    \loigiai{$T = 1/f = 1/50 = 0,02$ s.}
\end{ex}
%Môn: Hóa Học
\begin{ex}
    Tính pH của dung dịch có nồng độ ion $[H^+] = 10^{-3}$ M.
    \shortans{3}
    \loigiai{$pH = -\log[H^+] = -\log(10^{-3}) = 3$.}
\end{ex}
%Môn: Địa Lí
\begin{ex}
    Tính đến năm 2021, khu vực Đông Nam Á có tổng cộng bao nhiêu quốc gia?
    \shortans{11}
    \loigiai{Khu vực Đông Nam Á gồm 11 quốc gia: Việt Nam, Lào, Cam-pu-chia, Thái Lan, Mi-an-ma, Ma-lai-xi-a, Xin-ga-po, In-đô-nê-xi-a, Phi-líp-pin, Bru-nây và Đông Ti-mo.}
\end{ex}
%Môn: Lịch Sử
\begin{ex}
    Cuộc Duy tân Minh Trị ở Nhật Bản bắt đầu vào năm nào?
    \shortans{1868}
    \loigiai{Năm 1868, Thiên hoàng Minh Trị thực hiện cuộc cải cách đưa Nhật Bản phát triển theo con đường tư bản chủ nghĩa.}
\end{ex}
%Môn: Giáo Dục Kinh Tế và Pháp Luật
\begin{ex}
    Theo quy định, công dân đủ bao nhiêu tuổi trở lên có quyền ứng cử vào Quốc hội?
    \shortans{21}
    \loigiai{Công dân đủ 18 tuổi có quyền bầu cử, đủ 21 tuổi trở lên có quyền ứng cử.}
\end{ex}
%Môn: Toán
\begin{ex}
    Giá trị của $\cos(0^\circ)$ là bao nhiêu?
    \shortans{1}
    \loigiai{Dựa vào bảng giá trị lượng giác cơ bản, $\cos(0^\circ) = 1$.}
\end{ex}
%Môn: Vật Lí
\begin{ex}
    Một mạch điện có hiệu điện thế $U = 12$ V và cường độ dòng điện $I = 3$ A. Điện trở $R$ của mạch là bao nhiêu ohm?
    \shortans{4}
    \loigiai{$R = U/I = 12/3 = 4 \Omega$.}
\end{ex}
%Môn: Hóa Học
\begin{ex}
    Số oxi hóa của nguyên tử sulfur trong phân tử $SO_2$ là bao nhiêu?
    \shortans{4}
    \loigiai{Trong $SO_2$, $O$ là -2, tổng bằng 0 nên $S$ là +4.}
\end{ex}
%Môn: Địa Lí
\begin{ex}
    Theo bảng 9.1 trang 37, năm 1957 Cộng đồng kinh tế châu Âu có bao nhiêu nước thành viên?
    \shortans{6}
    \loigiai{Năm 1957, Cộng đồng kinh tế châu Âu được thành lập với 6 quốc gia thành viên.}
\end{ex}
%Môn: Tiếng Anh
\begin{ex}
    How many pillars does the ASEAN Community consist of?
    \shortans{3}
    \loigiai{The ASEAN Community is based on three pillars: Political-Security, Economic, and Socio-Cultural.}
\end{ex}
%Môn: Toán
\begin{ex}
    Tính đạo hàm của hàm số $y = 4x$ tại $x=1$.
    \shortans{4}
    \loigiai{$y' = 4$. Tại mọi $x$, giá trị đạo hàm luôn bằng 4.}
\end{ex}
%Môn: Vật Lí
\begin{ex}
    Một sóng cơ có tần số $f = 10$ Hz, tốc độ truyền sóng $v = 20$ m/s. Bước sóng $\lambda$ là bao nhiêu mét?
    \shortans{2}
    \loigiai{$\lambda = v/f = 20/10 = 2$ m.}
\end{ex}
%Môn: Hóa Học
\begin{ex}
    Phân tử ethanol ($C_2H_5OH$) có bao nhiêu nguyên tử carbon?
    \shortans{2}
    \loigiai{Công thức phân tử là $C_2H_6O$, chứa 2 nguyên tử C.}
\end{ex}
%Môn: Lịch Sử
\begin{ex}
    Liên Xô được thành lập vào năm nào?
    \shortans{1922}
    \loigiai{Tháng 12/1922, Liên bang Cộng hòa xã hội chủ nghĩa Xô viết (Liên Xô) chính thức ra đời.}
\end{ex}
%Môn: Giáo Dục Kinh Tế và Pháp Luật
\begin{ex}
    Tỉ lệ lạm phát của Việt Nam năm 2021 là bao nhiêu \%?
    \shortans{1,84}
    \loigiai{Theo biểu đồ trang 17, tỉ lệ lạm phát năm 2021 là 1,84\%.}
\end{ex}
%Môn: Toán
\begin{ex}
    Tính giới hạn $L = \lim_{x \to 2} (x + 3)$.
    \shortans{5}
    \loigiai{$L = 2 + 3 = 5$.}
\end{ex}
%Môn: Vật Lí
\begin{ex}
    Góc lệch cực đại của con lắc đơn trong dao động điều hòa phải nhỏ hơn hoặc bằng bao nhiêu độ?
    \shortans{10}
    \loigiai{Để con lắc đơn dao động điều hòa, góc lệch $\alpha_0 \le 10^\circ$.}
\end{ex}
%Môn: Hóa Học
\begin{ex}
    Sulfuric acid ($H_2SO_4$) có thể phân li tối đa bao nhiêu nấc proton $H^+$ trong nước?
    \shortans{2}
    \loigiai{Đây là acid 2 nấc: $H_2SO_4 \to H^+ + HSO_4^-$ và $HSO_4^- \to H^+ + SO_4^{2-}$.}
\end{ex}
%Môn: Địa Lí
\begin{ex}
    Năm 2020, tỉ lệ tăng dân số tự nhiên của khu vực Mỹ La-tinh là bao nhiêu \%?
    \shortans{0,94}
    \loigiai{Dựa vào hình 6.4 trang 26, tỉ lệ tăng dân số năm 2020 là 0,94\%.}
\end{ex}
%Môn: Lịch Sử
\begin{ex}
    Năm bao nhiêu thực dân Pháp bắt đầu nổ súng xâm lược Việt Nam tại Đà Nẵng?
    \shortans{1858}
    \loigiai{Ngày 1-9-1858, liên quân Pháp - Tây Ban Nha tấn công cửa biển Đà Nẵng.}
\end{ex}
%Môn: Toán
\begin{ex}
    Tìm nghiệm của phương trình $2^x = 16$.
    \shortans{4}
    \loigiai{$2^x = 2^4 \Rightarrow x = 4$.}
\end{ex}
%Môn: Vật Lí
\begin{ex}
    Một tụ điện có điện dung $C = 2$ $\mu F$ tích điện dưới hiệu điện thế $U = 100$ V. Điện tích $Q$ của tụ là bao nhiêu $\mu C$?
    \shortans{200}
    \loigiai{$Q = C \cdot U = 2 \cdot 100 = 200 \mu C$.}
\end{ex}
%Môn: Hóa Học
\begin{ex}
    Trong phân tử Alkyne có chứa ít nhất bao nhiêu liên kết ba $C \equiv C$?
    \shortans{1}
    \loigiai{Alkyne là hydrocarbon không no, mạch hở, trong phân tử có 1 liên kết ba.}
\end{ex}
%Môn: Địa Lí
\begin{ex}
    Kênh đào Pa-na-ma có chiều dài khoảng bao nhiêu km?
    \shortans{64}
    \loigiai{SGK Địa lí 11 trang 22 cung cấp thông tin chiều dài kênh là 64 km.}
\end{ex}
%Môn: Giáo Dục Kinh Tế và Pháp Luật
\begin{ex}
    Chọn đáp án Đúng hoặc Sai
    \choiceTF
    {\True Cung là khối lượng hàng hóa người bán sẵn sàng bán.}
    {Khi giá tăng, lượng cung thường giảm.}
    {\True Cầu là lượng hàng hóa người mua sẵn sàng mua tại mức giá nhất định.}
    {\True Giá cả thị trường là một nhân tố ảnh hưởng đến cung và cầu.}
    \loigiai{
        \begin{itemchoice}
        \itemch Đúng định nghĩa.
        \itemch giá tăng thì lượng cung tăng.
        \itemch Đúng định nghĩa.
        \itemch 
        \end{itemchoice}
    }
\end{ex}
%Môn: Toán
\begin{ex}
    Cho hình chóp tam giác đều $S.ABC$. Số mặt của hình chóp này là bao nhiêu?
    \shortans{4}
    \loigiai{Gồm 1 mặt đáy và 3 mặt bên.}
\end{ex}
%Môn: Vật Lí
\begin{ex}
    Tốc độ ánh sáng trong chân không xấp xỉ $3 \cdot 10^n$ m/s. Tìm số nguyên $n$.
    \shortans{8}
    \loigiai{$c \approx 300.000.000 = 3 \cdot 10^8$ m/s.}
\end{ex}
%Môn: Hóa Học
\begin{ex}
    Chất nào sau đây là hydrocarbon no?
    \choice
    {Ethylene}
    {Acetylene}
    {\True Ethane}
    {Benzene}
    \loigiai{Ethane ($C_2H_6$) chỉ chứa liên kết đơn nên là hydrocarbon no.}
\end{ex}
%Môn: Địa Lí
\begin{ex}
    Bra-xin có diện tích khoảng bao nhiêu triệu km$^2$?
    \shortans{8,5}
    \loigiai{Diện tích Bra-xin xấp xỉ 8,5 triệu km$^2$.}
\end{ex}
%Môn: Lịch Sử
\begin{ex}
    Sự kiện quần chúng nhân dân Pa-ri tấn công ngục Ba-xti diễn ra vào ngày 14 tháng mấy năm 1789?
    \shortans{7}
    \loigiai{Ngày 14-7-1789 là ngày quốc khánh của nước Pháp.}
\end{ex}
%Môn: Sinh Học
\begin{ex}
    Một chu kì tim ở người bình thường kéo dài khoảng bao nhiêu giây?
    \shortans{0,8}
    \loigiai{Chu kì tim gồm pha nhĩ co 0,1s, thất co 0,3s và dãn chung 0,4s.}
\end{ex}
%Môn: Toán
\begin{ex}
    Số hạng đầu $u_1$ của dãy số $u_n = 2n - 1$ là bao nhiêu?
    \shortans{1}
    \loigiai{$u_1 = 2(1) - 1 = 1$.}
\end{ex}
%Môn: Vật Lí
\begin{ex}
    Hiệu điện thế $U$ giữa hai bản tụ là 10V, điện tích $Q$ là 20C. Điện dung $C$ là bao nhiêu Fara?
    \shortans{2}
    \loigiai{$C = Q/U = 20/10 = 2$ F.}
\end{ex}
%Môn: Hóa Học
\begin{ex}
    Hợp chất $CH_3-CH_2-OH$ có tên gọi thay thế là gì?
    \choice
    {Methanol}
    {\True Ethanol}
    {Phenol}
    {Ethane}
    \loigiai{Hợp chất có 2 carbon và nhóm -OH nên là ethanol.}
\end{ex}
%Môn: Địa Lí
\begin{ex}
    \immini[thm]{
        Dựa vào biểu đồ 2.1 trang 17, hãy cho biết pH của nước biển là khoảng bao nhiêu?
    }{
        [Image 2.1]}
    \shortans{8}
    \loigiai{Nước biển có tính kiềm nhẹ, pH khoảng 8,0 - 8,4.}
\end{ex}
%Môn: Giáo Dục Kinh Tế và Pháp Luật
\begin{ex}
    Cạnh tranh lành mạnh mang lại lợi ích cho ai?
    \choice
    {Chỉ người sản xuất}
    {Chỉ người tiêu dùng}
    {Chỉ nhà nước}
    {\True Cả người sản xuất, người tiêu dùng và xã hội}
    \loigiai{Cạnh tranh lành mạnh thúc đẩy kinh tế phát triển và bảo vệ quyền lợi các bên.}
\end{ex}
%Môn: Toán
\begin{ex}
    Tính giá trị của $\log_3 9$.
    \shortans{2}
    \loigiai{$3^2 = 9 \Rightarrow \log_3 9 = 2$.}
\end{ex}
%Môn: Vật Lí
\begin{ex}
    Dao động có biên độ giảm dần theo thời gian gọi là dao động gì?
    \choice
    {Duy trì}
    {Cưỡng bức}
    {\True Tắt dần}
    {Điều hòa}
    \loigiai{Do tác dụng của lực ma sát, biên độ giảm dần nên gọi là dao động tắt dần.}
\end{ex}
%Môn: Hóa Học
\begin{ex}
    \sochc{2}{
        Xét các chất sau: $HCl, CH_3COOH, NaOH, C_{12}H_{22}O_{11}$.
    }
    \begin{chc}
        Có bao nhiêu chất là chất điện li mạnh?
        \shortans{2}
        \loigiai{$HCl$ và $NaOH$ là các chất điện li mạnh.}
    \end{chc}
    \begin{chc}
        Chất nào không dẫn điện khi tan trong nước?
        \loigiai{Saccharose ($C_{12}H_{22}O_{11}$) không phân li thành ion nên không dẫn điện.}
    \end{chc}
\end{ex}
%Môn: Địa Lí
\begin{ex}
    Thủ đô của nước Nga là thành phố nào?
    \choice
    {Xanh Pê-téc-bua}
    {\True Mát-xcơ-va}
    {Tô-ky-ô}
    {Oa-sinh-tơn}
    \loigiai{Mát-xcơ-va (Moscow) là trung tâm chính trị của Liên bang Nga.}
\end{ex}
%Môn: Lịch Sử
\begin{ex}
    Cuộc khởi nghĩa của Hai Bà Trưng diễn ra vào năm bao nhiêu sau Công nguyên?
    \shortans{40}
    \loigiai{Năm 40, Hai Bà Trưng phất cờ khởi nghĩa tại Hát Môn.}
\end{ex}
%Môn: Sinh Học
\begin{ex}
    Thực vật C4 có năng suất cao hơn thực vật C3 vì chúng không có hiện tượng nào?
    \choice
    {Quang hợp}
    {\True Hô hấp sáng}
    {Thoát hơi nước}
    {Hấp thụ khoáng}
    \loigiai{Hô hấp sáng làm lãng phí sản phẩm quang hợp ở thực vật C3.}
\end{ex}
%Môn: Toán
\begin{ex}
    Cho một khối lập phương. Số cạnh của nó là bao nhiêu?
    \shortans{12}
    \loigiai{Hình lập phương có 12 cạnh bằng nhau.}
\end{ex}
%Môn: Vật Lí
\begin{ex}
    Gia tốc trong dao động điều hòa biến thiên như thế nào so với li độ?
    \choice
    {Cùng pha}
    {\True Ngược pha}
    {Sớm pha $\pi/2$}
    {Trễ pha $\pi/2$}
    \loigiai{Biểu thức $a = -\omega^2 x$ cho thấy $a$ và $x$ luôn ngược pha.}
\end{ex}
%Môn: Hóa Học
\begin{ex}
    Phản ứng tráng bạc của aldehyde với thuốc thử Tollens tạo ra kim loại gì?
    \choice
    {Cu}
    {Au}
    {Fe}
    {\True Ag}
    \loigiai{Phản ứng khử ion bạc thành bạc kim loại ($Ag$).}
\end{ex}
%Môn: Địa Lí
\begin{ex}
    Nhật Bản nằm ở khu vực nào của châu Á?
    \choice
    {Đông Nam Á}
    {Tây Nam Á}
    {Trung Á}
    {\True Đông Á}
    \loigiai{Nhật Bản là quốc gia quần đảo nằm ở phía Đông của lục địa châu Á.}
\end{ex}
%Môn: Giáo Dục Kinh Tế và Pháp Luật
\begin{ex}
    Lạm phát vừa phải có mức tăng chỉ số giá tiêu dùng hàng năm là bao nhiêu \%?
    \choice
    {\True Dưới $10\%$}
    {Từ $10\%$ đến $100\%$}
    {Trên $1000\%$}
    {Bằng $0\%$}
    \loigiai{Lạm phát một con số (dưới $10\%$) được coi là lạm phát vừa phải.}
\end{ex}
%Môn: Toán
\begin{ex}
    Tính giới hạn $L = \lim_{n \to \infty} \frac{n+1}{n}$.
    \shortans{1}
    \loigiai{$L = \lim (1 + 1/n) = 1$.}
\end{ex}
%Môn: Vật Lí
\begin{ex}
    Tần số góc $\omega$ của con lắc lò xo được tính theo công thức nào?
    \choice
    {$\sqrt{g/l}$}
    {\True $\sqrt{k/m}$}
    {$2\pi \sqrt{m/k}$}
    {$m/k$}
    \loigiai{$\omega = \sqrt{k/m}$ trong đó $k$ là độ cứng, $m$ là khối lượng.}
\end{ex}
%Môn: Hóa Học
\begin{ex}
    Liên kết trong phân tử nitrogen ($N_2$) là liên kết gì?
    \choice
    {Đơn}
    {Đôi}
    {\True Ba}
    {Ion}
    \loigiai{Hai nguyên tử Nitrogen liên kết với nhau bằng liên kết ba bền vững $N \equiv N$.}
\end{ex}
%Môn: Địa Lí
\begin{ex}
    Thành phố nào được mệnh danh là siêu đô thị có dân số trên 10 triệu người ở Bra-xin?
    \choice
    {Ri-ô đê Gia-nê-rô}
    {\True Xao Pao-lô}
    {Bra-xi-li-a}
    {Li-ma}
    \loigiai{Xao Pao-lô là siêu đô thị lớn nhất Nam Mỹ.}
\end{ex}
%Môn: Tiếng Anh
\begin{ex}
    The prefix 'self-' in 'self-study' means ______.
    \choice
    {With others}
    {\True By oneself}
    {With a teacher}
    {Quickly}
    \loigiai{Self-study is learning by yourself without formal instruction.}
\end{ex}
%Môn: Lịch Sử
\begin{ex}
    Chiến thắng Bạch Đằng năm 938 do Ngô Quyền lãnh đạo đánh bại quân xâm lược nào?
    \choice
    {Nam Hán}
    {\True Nam Hán}
    {Tống}
    {Mông Cổ}
    \loigiai{Ngô Quyền đánh tan quân Nam Hán trên sông Bạch Đằng.}
\end{ex}
%Môn: Toán
\begin{ex}
    Cho $A, B$ là hai biến cố độc lập. Biết $P(A)=0,5$ và $P(B)=0,6$. Tính $P(AB)$.
    \shortans{0,3}
    \loigiai{$P(AB) = 0,5 \cdot 0,6 = 0,3$.}
\end{ex}
%Môn: Sinh Học
\begin{ex}
    Quá trình quang hợp giải phóng khí gì ra môi trường?
    \loigiai{Quang hợp hấp thụ $CO_2$ và giải phóng $O_2$.}
\end{ex}
%Môn: Hóa Học
\begin{ex}
    Chọn đáp án đúng hoặc sai
    \choiceTF
    {\True Methane là thành phần chính của khí mỏ dầu.}
    {Alkane có phản ứng cộng với nước bromine.}
    {\True Phản ứng đặc trưng của alkane là phản ứng thế.}
    {Alkane tan tốt trong nước.}
    \loigiai{
        \begin{itemchoice}
        \itemch 
        \itemch alkane không có liên kết bội.
        \itemch 
        \itemch hydrocarbon không phân cực nên không tan trong nước.
        \end{itemchoice}
    }
\end{ex}
%Môn: Vật Lí
\begin{ex}
    Đơn vị của cường độ dòng điện là gì?
    \choice
    {Vôn}
    {Oát}
    {\True Ampe}
    {Ôm}
    \loigiai{Ampe (A) là đơn vị đo cường độ dòng điện.}
\end{ex}
%Môn: Địa Lí
\begin{ex}
    Khu vực nào giàu tài nguyên dầu mỏ và khí tự nhiên nhất thế giới?
    \choice
    {Đông Nam Á}
    {\True Tây Nam Á}
    {Mỹ La-tinh}
    {Châu Âu}
    \loigiai{Tây Nam Á sở hữu trữ lượng dầu mỏ khổng lồ, đặc biệt quanh vịnh Péc-xích.}
\end{ex}
%Môn: Toán
\begin{ex}
    Tìm giá trị của $\sin(90^\circ)$.
    \shortans{1}
    \loigiai{$\sin(90^\circ) = 1$.}
\end{ex}
%Môn: Sinh Học
\begin{ex}
    Hệ thống ống khí là cơ quan hô hấp của nhóm động vật nào?
    \choice
    {Thú}
    {Chim}
    {\True Côn trùng}
    {Cá}
    \loigiai{Côn trùng trao đổi khí trực tiếp qua hệ thống ống khí.}
\end{ex}
%Môn: Hóa Học
\begin{ex}
    Dung dịch ammonia ($NH_3$) làm quỳ tím chuyển sang màu gì?
    \choice
    {Đỏ}
    {\True Xanh}
    {Vàng}
    {Không đổi màu}
    \loigiai{Ammonia có tính base yếu.}
\end{ex}
%Môn: Địa Lí
\begin{ex}
    Siêu đô thị nào sau đây nằm ở Hoa Kỳ?
    \choice
    {Luân Đôn}
    {Pa-ri}
    {\True Niu Oóc}
    {Bắc Kinh}
    \loigiai{Niu Oóc (New York) là trung tâm kinh tế lớn nhất Hoa Kỳ.}
\end{ex}
%Môn: Lịch Sử
\begin{ex}
    Cuộc khởi nghĩa Lam Sơn kết thúc thắng lợi vào năm nào?
    \shortans{1427}
    \loigiai{Năm 1427, nghĩa quân Lam Sơn quét sạch quân Minh ra khỏi bờ cõi.}
\end{ex}
%Môn: Toán
\begin{ex}
    Tính đạo hàm của hàm số $y = x^2$.
    \loigiai{$y' = 2x$.}
\end{ex}
%Môn: Vật Lí
\begin{ex}
    Sóng âm có tần số dưới 16 Hz được gọi là gì?
    \choice
    {Âm nghe được}
    {Siêu âm}
    {\True Hạ âm}
    {Nhạc âm}
    \loigiai{Dưới 16 Hz tai người không nghe được, gọi là hạ âm.}
\end{ex}
%Môn: Giáo Dục Kinh Tế và Pháp Luật
\begin{ex}
    Hành vi nào vi phạm quyền bất khả xâm phạm về thân thể?
    \choice
    {Khen thưởng người tốt}
    {Giúp đỡ người già}
    {\True Bắt giữ người trái pháp luật}
    {Tuyên truyền pháp luật}
    \loigiai{Tự ý bắt giữ người khi không có lệnh là vi phạm pháp luật.}
\end{ex}
%Môn: Địa Lí
\begin{ex}
    Thành viên thứ 11 vừa được ASEAN đồng ý về mặt nguyên tắc năm 2022 là nước nào?
    \choice
    {Bru-nây}
    {Lào}
    {\True Đông Ti-mo}
    {Việt Nam}
    \loigiai{Đông Ti-mo (Timor-Leste) là thành viên tiềm năng thứ 11.}
\end{ex}
%Môn: Hóa Học
\begin{ex}
    Liên kết $O-H$ trong phân tử alcohol là liên kết gì?
    \choice
    {Phi cực}
    {\True Phân cực}
    {Ion}
    {Kim loại}
    \loigiai{Do độ âm điện của $O$ lớn hơn $H$ nên liên kết bị phân cực về phía $O$.}
\end{ex}
%Môn: Toán
\begin{ex}
    Dãy số $1, 1, 1, 1, ...$ là cấp số cộng có công sai $d$ bằng bao nhiêu?
    \shortans{0}
    \loigiai{$d = 1 - 1 = 0$.}
\end{ex}
%Môn: Vật Lí
\begin{ex}
    Trong dao động điều hòa, li độ $x$ đạt cực đại tại đâu?
    \choice
    {Vị trí cân bằng}
    {\True Vị trí biên}
    {Vị trí bất kì}
    {Không có cực đại}
    \loigiai{Tại biên, $|x| = A$.}
\end{ex}
%Môn: Lịch Sử
\begin{ex}
    Vua nào đã ban chiếu dời đô về Thăng Long năm 1010?
    \choice
    {Lý Thái Tông}
    {\True Lý Thái Tổ}
    {Lý Nhân Tông}
    {Lý Thánh Tông}
    \loigiai{Lý Công Uẩn (Lý Thái Tổ) dời đô từ Hoa Lư về Đại La.}
\end{ex}
%Môn: Sinh Học
\begin{ex}
    Phổi chim có cấu tạo đặc biệt với hệ thống gì giúp hô hấp hiệu quả?
    \choice
    {Phế nang}
    {\True Ống khí}
    {Túi khí}
    {Mang}
    \loigiai{Chim hô hấp qua các ống khí, không có phế nang.}
\end{ex}
%Môn: Địa Lí
\begin{ex}
    Châu lục nào có diện tích nhỏ nhất thế giới?
    \loigiai{Châu Đại Dương (O-IA).}
\end{ex}
%Môn: Hóa Học
\begin{ex}
    Sản phẩm của phản ứng giữa $CH_3COOH$ và $NaOH$ là muối gì?
    \loigiai{Sodium acetate ($CH_3COONa$).}
\end{ex}
%Môn: Toán
\begin{ex}
    Cho góc $\alpha$ thuộc góc phần tư thứ II. Dấu của $\sin \alpha$ là gì?
    \loigiai{Ở góc phần tư thứ II, tung độ dương nên $\sin \alpha > 0$.}
\end{ex}
%Môn: Vật Lí
\begin{ex}
    Chọn đáp án Đúng hoặc Sai.
    \choiceTF
    {\True Sóng ngang truyền được trong chất rắn.}
    {Sóng dọc không truyền được trong chất lỏng.}
    {\True Sóng âm là sóng dọc trong chất khí.}
    {Vận tốc âm thanh trong thép nhỏ hơn trong nước.}
    \loigiai{
        \begin{itemchoice}
        \itemch 
        \itemch Sai.
        \itemch 
        \itemch âm truyền nhanh nhất trong chất rắn.
        \end{itemchoice}
    }
\end{ex}
%Môn: Địa Lí
\begin{ex}
    Australia (Ô-xtrây-li-a) là quốc gia chiếm cả một ______.
    \loigiai{Lục địa (L-D).}
\end{ex}
%Môn: Tiếng Anh
\begin{ex}
    'Generation gap' usually causes ______ between parents and children.
    \choice
    {Agreement}
    {Happiness}
    {\True Conflicts}
    {Cooperation}
    \loigiai{Differences in viewpoints lead to arguments or conflicts.}
\end{ex}
%Môn: Toán
\begin{ex}
    Số nghiệm của phương trình $\sin x = 2$ là bao nhiêu?
    \shortans{0}
    \loigiai{Vì $-1 \le \sin x \le 1$ nên phương trình vô nghiệm.}
\end{ex}
%Môn: Hóa Học
\begin{ex}
    Tính chất hóa học đặc trưng của hợp chất carbonyl ($C=O$) là phản ứng gì?
    \choice
    {Thế}
    {\True Cộng}
    {Tách}
    {Thủy phân}
    \loigiai{Nhóm $C=O$ dễ tham gia phản ứng cộng vào liên kết đôi.}
\end{ex}
%Môn: Vật Lí
\begin{ex}
    Trong dao động tắt dần, đại lượng nào giảm dần theo thời gian?
    \choice
    {Tần số}
    {Chu kì}
    {\True Biên độ}
    {Pha ban đầu}
    \loigiai{Năng lượng và biên độ giảm do ma sát.}
\end{ex}
%Môn: Sinh Học
\begin{ex}
    Sự kết hợp giữa giao tử đực và giao tử cái tạo thành gì?
    \loigiai{Hợp tử (H-T).}
\end{ex}
%Môn: Lịch Sử
\begin{ex}
    Cách mạng tháng Tám ở Việt Nam diễn ra vào năm nào?
    \shortans{1945}
    \loigiai{Cuộc cách mạng giành độc lập năm 1945.}
\end{ex}
%Môn: Địa Lí
\begin{ex}
    Dân số Hoa Kỳ năm 2020 là khoảng bao nhiêu triệu người?
    \shortans{331}
    \loigiai{Dân số xấp xỉ 331,5 triệu người.}
\end{ex}
%Môn: Toán
\begin{ex}
    Khoảng cách từ điểm $M(0;0;1)$ đến mặt phẳng $z=0$ là bao nhiêu?
    \shortans{1}
    \loigiai{$d = |1-0| = 1$.}
\end{ex}
%Môn: Hóa Học
\begin{ex}
    Hợp chất nào có mùi thơm của chuối chín?
    \loigiai{Isoamyl acetate (IA).}
\end{ex}
%Môn: Giáo Dục Kinh Tế và Pháp Luật
\begin{ex}
    Trong kinh tế, 'Cầu' biểu thị lượng hàng hóa mà người mua ______.
    \choice
    {Cần mua}
    {Thích mua}
    {\True Sẵn sàng mua và có khả năng thanh toán}
    {Đã mua xong}
    \loigiai{Cầu phải đi kèm khả năng chi trả.}
\end{ex}
%Môn: Địa Lí
\begin{ex}
    Quốc gia nào sau đây không thuộc EU?
    \choice
    {Đức}
    {Pháp}
    {Ý}
    {\True Việt Nam}
    \loigiai{Việt Nam thuộc ASEAN.}
\end{ex}
%Môn: Tiếng Anh
\begin{ex}
    A place where students receive higher education is called a ______.
    \loigiai{University (U).}
\end{ex}
%Môn: Toán
\begin{ex}
    Nếu $\sin \alpha = 0,6$ thì $\sin^2 \alpha$ bằng bao nhiêu?
    \shortans{0,36}
    \loigiai{$0,6 \cdot 0,6 = 0,36$.}
\end{ex}
%Môn: Vật Lí
\begin{ex}
    Độ cao của âm phụ thuộc vào đại lượng nào?
    \choice
    {Biên độ}
    {Cường độ}
    {\True Tần số}
    {Đồ thị dao động}
    \loigiai{Tần số càng lớn âm càng cao.}
\end{ex}
%Môn: Hóa Học
\begin{ex}
    Phenol là hợp chất có nhóm $-OH$ liên kết trực tiếp với nguyên tử carbon của ______.
    \loigiai{Vòng Benzene (V-B).}
\end{ex}
%Môn: Sinh Học
\begin{ex}
    Cơ quan nào ở người làm nhiệm vụ lọc máu tạo nước tiểu?
    \loigiai{Thận (T).}
\end{ex}
%Môn: Lịch Sử
\begin{ex}
    Vua cuối cùng của chế độ phong kiến Việt Nam là ai?
    \loigiai{Vua Bảo Đại (B-D).}
\end{ex}
%Môn: Địa Lí
\begin{ex}
    Úc (Australia) nằm ở bán cầu nào?
    \choice
    {Bắc}
    {\True Nam}
    {Tây}
    {Đông}
    \loigiai{Toàn bộ lục địa Úc nằm ở bán cầu Nam.}
\end{ex}
%Môn: Toán
\begin{ex}
    Số cách xếp 3 người vào một hàng dọc là $3!$. Kết quả là bao nhiêu?
    \shortans{6}
    \loigiai{$3 \cdot 2 \cdot 1 = 6$.}
\end{ex}
%Môn: Vật Lí
\begin{ex}
    Tốc độ truyền sóng trong môi trường nào là nhanh nhất?
    \choice
    {Lỏng}
    {Khí}
    {\True Rắn}
    {Chân không}
    \loigiai{Mật độ phân tử càng lớn tốc độ truyền càng nhanh.}
\end{ex}
%Môn: Hóa Học
\begin{ex}
    Tính khối lượng mol của $H_2SO_4$. (H=1, S=32, O=16)
    \shortans{98}
    \loigiai{$2\cdot 1 + 32 + 16\cdot 4 = 98$.}
\end{ex}
%Môn: Giáo Dục Kinh Tế và Pháp Luật
\begin{ex}
    Thất nghiệp do cơ cấu xảy ra khi có sự thay đổi về ______.
    \loigiai{Công nghệ (C-N) hoặc cấu trúc kinh tế.}
\end{ex}
